\documentclass[12pt]{article}

\usepackage[a4paper, margin=2.5cm]{geometry}
\usepackage{amsmath}
\usepackage{amssymb}
\usepackage{amsthm}
\usepackage[colorlinks=true, linkcolor=blue, citecolor=blue, urlcolor=blue]{hyperref}
\usepackage{booktabs}
\usepackage{natbib}

\newtheorem{theorem}{Theorem}
\newtheorem{proposition}[theorem]{Proposition}
\newtheorem{definition}[theorem]{Definition}
\newtheorem{remark}[theorem]{Remark}

\newcommand{\Mcl}{\mathcal{M}_{\mathrm{cl}}}
\newcommand{\Ccal}{\mathcal{C}}
\newcommand{\Fcal}{\mathcal{F}}
\newcommand{\Lcal}{\mathcal{L}}
\newcommand{\Sadm}{\mathcal{S}_{\mathrm{adm}}}
\newcommand{\Scl}{\mathcal{S}_{\mathrm{cl}}}
\newcommand{\phig}{\varphi}

\title{Closure Dynamics:\\A Variational and Gradient Formulation\\of Constraint-Manifold Evolution}
\author{%
  Lee Smart\\[4pt]
  \textit{Institute of Vibrational Field Dynamics}\\[2pt]
  \texttt{contact@vibrationalfielddynamics.org}\\[2pt]
  \texttt{@vfd\_org}%
}
\date{April 2026}

\begin{document}

\maketitle

\noindent\textbf{Status:} Preprint (not peer-reviewed)

\begin{abstract}
Papers~V--XI developed a structural constraint framework for particle mass, gauge structure, confinement, continuous geometry, and a gravitational analogy, but did not include a time-evolution law. The present paper introduces minimal dynamics into the closure framework by defining state evolution as gradient flow of the closure functional $\Fcal$.

Under this evolution, the closure functional decreases monotonically along trajectories, closure-stable states ($\Fcal = 0$) are fixed points, and distinct attractor basins correspond to distinct particle classes. Disconnected configurations are dynamically unstable: they flow toward connected composites, providing a dynamical interpretation of confinement. A conservative variational extension (Lagrangian formulation) is introduced as a second-order generalisation.

The evolution parameter $t$ is not identified with physical time. No relativistic invariance, quantum field theory, scattering amplitudes, or coupling constants are derived. The paper establishes the minimal dynamical layer consistent with the closure-constraint structure, providing the formal transition from a purely structural framework to one that admits evolution.
\end{abstract}

%==================================================================
\section{Introduction}\label{sec:intro}
%==================================================================

The preceding papers in this programme~\citep{SmartV, SmartVII, SmartVIII, SmartIX, SmartX, SmartXI} developed a constraint-based framework in which physical states correspond to the closed state set $\Scl$, mass is the closure residual, and the constraint manifold $\Mcl$ carries geometric structure. Throughout that development, the framework was purely structural: it specified \textit{which} states are admissible and \textit{what} properties they carry, but not \textit{how} states evolve.

The present paper fills this gap by introducing a minimal evolution law. The dynamics is defined as gradient flow of the closure functional: states evolve in the direction that most rapidly reduces constraint violation. This is the simplest dynamical extension consistent with the closure structure, and it produces the expected behaviour: closure-stable states are attractors, disconnected configurations are unstable, and the constraint manifold governs the geometry of admissible motion.

\subsection*{Scope and epistemic status}

This paper introduces \textit{minimal dynamics}: a gradient-flow evolution law and its variational extension. It does not introduce:
\begin{itemize}
    \item physical time (the evolution parameter $t$ is auxiliary),
    \item relativistic invariance or Lorentz covariance,
    \item quantum field theory or operator-valued dynamics,
    \item scattering amplitudes, decay rates, or coupling constants.
\end{itemize}

The dynamics is presented as the simplest evolution law compatible with the closure structure. Whether it can be extended to a full physical dynamics is an open question.

%==================================================================
\section{Time-Dependent State Space}\label{sec:timedep}
%==================================================================

We extend the admissible state space to include a time parameter:
\begin{equation}\label{eq:timedep}
    \psi = \psi(t), \qquad t \in \mathbb{R}
\end{equation}
where $\psi(t) \in \Sadm$ for all $t$. The state traces a curve in the admissible domain as the evolution parameter advances.

\begin{remark}[Status of $t$]
The parameter $t$ is introduced as an auxiliary evolution parameter governing relaxation within the constraint framework. It is not identified with physical coordinate time, proper time, or any relativistic time parameter. The identification of $t$ with a physical quantity would require additional structure (e.g.\ a metric signature or Lorentz group action) not present in this paper.
\end{remark}

%==================================================================
\section{Gradient Flow Dynamics}\label{sec:gradflow}
%==================================================================

\subsection{The evolution law}

We define the evolution of states as gradient flow of the closure functional $\Fcal$:
\begin{equation}\label{eq:gradientflow}
    \frac{d\psi}{dt} = -\nabla \Fcal(\psi)
\end{equation}
where $\nabla \Fcal$ is the gradient of $\Fcal$ with respect to the inner product on the state space $V$.

\subsection{Constrained evolution}

Since $\psi(t)$ must remain in $\Sadm$ (e.g.\ normalised), the flow is projected onto the admissible domain:
\begin{equation}\label{eq:projectedflow}
    \frac{d\psi}{dt} = -\Pi_{\Sadm}\!\left(\nabla \Fcal(\psi)\right)
\end{equation}
where $\Pi_{\Sadm}$ denotes projection onto the tangent space of $\Sadm$ at $\psi$.

\subsection{Monotonic decrease}

\begin{proposition}[Monotonic descent]\label{prop:descent}
Along trajectories of the projected gradient flow~\eqref{eq:projectedflow}, the closure functional decreases monotonically:
\begin{equation}
    \frac{d\Fcal}{dt} = \langle \nabla\Fcal, \dot{\psi} \rangle = -\|\Pi_{\Sadm}(\nabla\Fcal)\|^2 \leq 0
\end{equation}
Equality holds if and only if $\Pi_{\Sadm}(\nabla\Fcal(\psi)) = 0$, i.e.\ at a constrained critical point.
\end{proposition}

\begin{proof}
By the chain rule, $d\Fcal/dt = \langle\nabla\Fcal, d\psi/dt\rangle$. Substituting~\eqref{eq:projectedflow}: $d\Fcal/dt = -\langle\nabla\Fcal, \Pi_{\Sadm}(\nabla\Fcal)\rangle$. Since $\Pi_{\Sadm}$ is an orthogonal projection, $\langle\nabla\Fcal, \Pi_{\Sadm}(\nabla\Fcal)\rangle = \|\Pi_{\Sadm}(\nabla\Fcal)\|^2 \geq 0$.
\end{proof}

\subsection{Interpretation}

The gradient flow drives the system toward lower values of $\Fcal$:
\begin{itemize}
    \item $\Fcal(\psi)$ decreases monotonically along trajectories.
    \item States with $\Fcal(\psi) = 0$ --- i.e.\ states in $\Scl$ --- are fixed points.
    \item The flow represents constraint relaxation: the system evolves in the direction of steepest descent toward closure satisfaction.
\end{itemize}

This is the simplest dynamics consistent with the closure structure: it uses only $\Fcal$ and the inner product, introduces no additional parameters, and automatically identifies $\Scl$ as the set of attractors. The gradient flow should be interpreted as a coarse-grained or relaxation dynamics, not a fundamental microscopic evolution law.

%==================================================================
\section{Fixed Points and Stability}\label{sec:fixedpoints}
%==================================================================

\begin{definition}[Dynamical fixed point]\label{def:fixedpoint}
A state $\psi^* \in \Sadm$ is a fixed point of the gradient-flow dynamics if:
\begin{equation}
    \Pi_{\Sadm}\!\left(\nabla\Fcal(\psi^*)\right) = 0
\end{equation}
\end{definition}

All states in $\Scl$ satisfy $\Fcal(\psi^*) = 0$ and are therefore fixed points: the gradient vanishes at the global minimum of $\Fcal$ on $\Sadm$.

\subsection{Local stability}

The stability of a fixed point is determined by the constrained Hessian:
\begin{equation}\label{eq:hessian}
    H(\psi) = \Pi_{\Sadm}\!\left(\nabla^2 \Fcal(\psi)\right)\big|_{T_\psi \Sadm}
\end{equation}
restricted to the tangent space of $\Sadm$ at $\psi$.

\begin{remark}[Stability criterion]
A fixed point $\psi^*$ is a stable attractor if $H(\psi^*)$ is positive semi-definite on $T_{\psi^*} \Sadm$: all nearby perturbations within the admissible domain are driven back toward $\psi^*$ by the flow. Saddle points of $\Fcal$ (with indefinite Hessian) are unstable fixed points.
\end{remark}

%==================================================================
\section{Attractor Structure}\label{sec:attractors}
%==================================================================

\begin{definition}[Attractor basin]\label{def:basin}
The attractor basin of a stable fixed point $\psi^* \in \Scl$ is:
\begin{equation}
    B(\psi^*) = \left\{\psi_0 \in \Sadm \;\middle|\; \lim_{t \to \infty} \psi(t; \psi_0) = \psi^*\right\}
\end{equation}
where $\psi(t; \psi_0)$ is the trajectory of the gradient flow starting from $\psi_0$.
\end{definition}

Within the framework, distinct attractor basins provide the dynamical analogue of distinct particle classes: initial configurations in basin $B(\psi^*)$ relax toward the closure-stable state $\psi^*$. The particle type is determined by which attractor the initial configuration flows toward.

\begin{remark}
The attractor-basin structure depends on the global geometry of $\Fcal$ on $\Sadm$. In general, the basin boundaries are the stable manifolds of saddle points. A full characterisation of the basin structure for the 600-cell closure functional is beyond the scope of the present paper.
\end{remark}

%==================================================================
\section{Mass and Dynamics}\label{sec:mass}
%==================================================================

Mass is defined as the closure residual (Papers~V, VII):
\begin{equation}
    m = \kappa\, d(\psi, \Mcl)
\end{equation}

Under gradient flow, this acquires a dynamical interpretation:

\begin{itemize}
    \item \textbf{Relaxation toward closure:} states with $\Fcal > 0$ evolve toward $\Mcl$. The closure residual $\Delta(\psi) = d(\psi, \Mcl)$ decreases along trajectories (since $\Fcal$ decreases and $\Delta$ is related to $\Fcal$).
    \item \textbf{Mass and the closure landscape:} mass is associated with the local structure of the closure landscape around $\Mcl$, rather than requiring states to lie permanently off the manifold in equilibrium. Physical particles correspond to closure-stable configurations; their mass reflects the closure-residual structure at those configurations.
    \item \textbf{Motion on $\Mcl$:} once a trajectory reaches $\Mcl$ (where $\Fcal = 0$), the gradient vanishes and the state is stationary under the dissipative gradient flow. Motion along $\Mcl$ requires the conservative (Lagrangian) extension of Section~\ref{sec:lagrangian}.
\end{itemize}

%==================================================================
\section{Confinement in Dynamics}\label{sec:confinement}
%==================================================================

Disconnected configurations (Paper~VIII) have $\Fcal > 0$ due to the gap penalty $\Xi > 0$. Under gradient flow:

\begin{itemize}
    \item disconnected states are \textit{not} fixed points (they have $\nabla\Fcal \neq 0$),
    \item they flow toward configurations with lower $\Fcal$,
    \item the flow drives them toward connected composites (where $\Xi = 0$),
    \item they cannot stabilise independently.
\end{itemize}

This provides a \textit{dynamical interpretation of confinement}: disconnected quark-like states are unstable under closure dynamics. They are not held together by a force; they are driven toward connected composites by the gradient of the constraint functional.

\begin{remark}
This dynamical confinement is a consequence of the gradient-flow structure and the gap penalty in $\Fcal$. It does not invoke a non-Abelian gauge field, a confining potential, or colour flux tubes. It is a structural result within the closure framework.
\end{remark}

%==================================================================
\section{Variational / Lagrangian Extension}\label{sec:lagrangian}
%==================================================================

The gradient flow~\eqref{eq:gradientflow} is a first-order, dissipative dynamics: it drives the system toward fixed points but does not describe oscillatory or conservative motion. To extend the framework to conservative dynamics, we introduce a Lagrangian formulation.

\subsection{The closure Lagrangian}

\begin{definition}[Closure Lagrangian]\label{def:lagrangian}
The closure Lagrangian is:
\begin{equation}\label{eq:lagrangian}
    \Lcal(\psi, \dot{\psi}) = \frac{1}{2}\|\dot{\psi}\|^2 - \Fcal(\psi)
\end{equation}
where $\|\dot{\psi}\|^2 = \langle\dot{\psi}, \dot{\psi}\rangle$ is the kinetic term and $\Fcal(\psi)$ plays the role of the potential. This choice corresponds to the simplest quadratic kinetic term compatible with the ambient inner-product structure; it is not claimed to be unique.
\end{definition}

\subsection{The action and equations of motion}

The action is:
\begin{equation}\label{eq:action}
    S[\psi] = \int_{t_1}^{t_2} \Lcal(\psi, \dot{\psi})\, dt
\end{equation}

Stationarity of the action ($\delta S = 0$) yields the Euler--Lagrange equation:
\begin{equation}\label{eq:EL}
    \ddot{\psi} = -\nabla\Fcal(\psi)
\end{equation}

This is a second-order dynamical system in which $\Fcal$ acts as a potential energy landscape. States oscillate around minima of $\Fcal$ (closure-stable configurations) rather than dissipating toward them.

\subsection{Constrained Lagrangian dynamics}

For evolution constrained to $\Sadm$, the equations of motion become:
\begin{equation}\label{eq:constrainedEL}
    \Pi_{\Sadm}\!\left(\ddot{\psi}\right) = -\Pi_{\Sadm}\!\left(\nabla\Fcal(\psi)\right)
\end{equation}
with the constraint that $\psi(t) \in \Sadm$ enforced throughout.

\subsection{Interpretation}

\begin{itemize}
    \item \textbf{Gradient flow} (first-order): dissipative relaxation toward closure. Models constraint satisfaction.
    \item \textbf{Lagrangian dynamics} (second-order): conservative motion on the $\Fcal$ landscape. Models evolution within and around closure-stable configurations.
\end{itemize}

Both formulations share the same fixed points ($\Scl$) and the same stability structure (determined by $H(\psi)$). They differ in their dynamical character: the gradient flow is irreversible; the Lagrangian dynamics is time-reversible.

\begin{remark}
The Lagrangian~\eqref{eq:lagrangian} is the simplest conservative extension of the closure structure. It uses only $\Fcal$ and the inner product, introduces no additional parameters, and reduces to the gradient flow in the overdamped limit ($\ddot{\psi} \to 0$). Whether a more physical Lagrangian (incorporating, e.g., gauge coupling or relativistic structure) can be constructed from the closure framework is an open question.
\end{remark}

%==================================================================
\section{Two Modes of Motion}\label{sec:modes}
%==================================================================

The dynamics defined in this paper produces two distinct modes of motion:

\begin{enumerate}
    \item \textbf{Transverse motion} (toward/away from $\Mcl$): evolution in the direction normal to the constraint manifold. Under gradient flow, this is relaxation toward closure. Under the Lagrangian, it is oscillation around $\Mcl$. This mode is associated with mass ($\Delta > 0$) and closure residual.

    \item \textbf{Tangential motion} (along $\Mcl$): evolution within the constraint manifold. Under the Lagrangian, states on $\Mcl$ can move along geodesic-like paths (Paper~X~\citep{SmartX}). Under gradient flow, tangential motion does not occur (the gradient vanishes on $\Mcl$).
\end{enumerate}

The separation into transverse and tangential modes provides a dynamical counterpart to the structural distinction between mass (distance from $\Mcl$) and geometry (motion on $\Mcl$).

\begin{remark}[Two levels of dynamics]
The gradient-flow dynamics should be interpreted as a relaxation process toward closure-stable configurations. It is not claimed to represent fundamental microscopic dynamics. A fully physical evolution law would require additional structure --- potentially building on the variational formulation, incorporating gauge coupling or relativistic invariance --- and represents the primary open problem for the programme's dynamical layer.
\end{remark}

%==================================================================
\section{Limitations}\label{sec:limits}
%==================================================================

\begin{itemize}
    \item \textbf{The evolution parameter $t$ is not physical time.} It is an auxiliary parameter. Identification with physical time requires additional structure (metric signature, Lorentz covariance) not present here.
    \item \textbf{No relativistic invariance.} The dynamics is defined on a state space with Euclidean inner product, not on a Lorentzian manifold.
    \item \textbf{No quantum dynamics.} The evolution is classical (deterministic gradient flow or Lagrangian mechanics). Quantisation is not addressed.
    \item \textbf{No scattering or decay.} The dynamics describes relaxation and oscillation around attractors, not scattering processes between states.
    \item \textbf{No coupling constants derived.} The mass parameter $\kappa$ and the closure-functional weights remain structural parameters.
    \item \textbf{The Lagrangian is minimal.} It is the simplest conservative extension; more physical Lagrangians may be needed.
    \item \textbf{Basin structure not fully characterised.} The global attractor-basin geometry depends on $\Fcal$ and is not computed here.
\end{itemize}

%==================================================================
\section{Conclusion}\label{sec:conclusion}
%==================================================================

The closure framework, previously purely structural, now admits a minimal dynamical extension. States evolve via gradient flow of the closure functional: $\Fcal$ decreases monotonically, closure-stable states are attractors, and disconnected configurations are dynamically unstable (confinement). A conservative Lagrangian extension provides second-order dynamics with oscillatory motion around closure-stable configurations.

The dynamics separates into two modes: transverse (toward/away from $\Mcl$, associated with mass) and tangential (along $\Mcl$, associated with geometry). This decomposition connects the dynamical evolution to the structural properties established in Papers~V--X.

The evolution parameter $t$ is not identified with physical time, and no relativistic or quantum structure is introduced. The paper establishes the formal transition from a constraint-based structural framework to one that admits evolution, providing the foundation on which more physical dynamics (gauge coupling, quantisation, scattering) could in principle be built.

%==================================================================
\bibliographystyle{plainnat}
\bibliography{references}

\end{document}
